% SHARD-ID: CA-26-GAUSSIAN-BOSON-GENERATOR-ALGEBRA
% SHARD-TITLE: One-Particle Generator Algebra for Gaussian Bosons
% SHARD-SUMMARY: Defines the one-particle H, P, J, and K candidates on a smooth momentum-space core.
% SHARD-SUMMARY: Derives the commutators that are exact for any scalar dispersion and isolates the residuals that depend on omega.
% SHARD-SUMMARY: Explains how second quantization transports the algebraic checks to the many-body Gaussian Fock space.
% SHARD-KEYWORDS: Gaussian boson, Poincare generators, boost, rotation, commutator, second quantization

\section{One-Particle Generator Algebra for Gaussian Bosons}
\label{sec:gaussian-boson-generator-algebra}

\subsection{Status, Sources, and Dependencies}

\textbf{Status: checked local derivation on a common core.}  The Poincare
target, Stone convention, and free-boson Fock realization are sourced.  The
commutator identities in this shard are local differential-operator derivations
on a smooth momentum-space core.

Dependencies: CA-11, CA-24--CA-25, and CONVENTIONS.md (h), (j).

% Source: references/cft/Schottenloher2008/Schottenloher2008.md:4040 -- Stone convention for one-parameter unitary groups.
% Source: references/cft/Schottenloher2008/Schottenloher2008.md:4092 -- Poincare group and relativistic invariance.
% Source: references/cft/Schottenloher2008/Schottenloher2008.md:4106 -- unitary Poincare representation.
% Source: references/cft/Schottenloher2008/Schottenloher2008.md:4113 -- translation generators P_mu and P_0=H.
% Source: references/cft/Schottenloher2008/Schottenloher2008.md:4244 -- free bosonic QFT carries a natural Poincare action on Fock space.
% Checked: test/runtests.jl, "Gaussian boson boost-time symbols".

\subsection{Momentum Core}

Work on a connected smooth momentum patch \(\mathcal U\) where
\(\omega(k)>0\).  Let
\[
  \mathcal D_0=C_c^\infty(\mathcal U)
\]
or the corresponding trigonometric-polynomial core on a finite periodic
Brillouin grid.  Define
\[
  X_a=i\partial_{k_a},
  \qquad
  h=\omega(k),
  \qquad
  p_a=k_a .
\]
The one-particle Hamiltonian and momenta are multiplication operators:
\[
  H^{(1)}=h,\qquad P_a^{(1)}=p_a.
\]
The rotation and time-zero boost candidates are
\[
  J_{ab}^{(1)}=X_aP_b-X_bP_a,
  \qquad
  K_a^{(1)}=\frac{1}{2}(X_ah+hX_a).
\]
All commutators below are identities on \(\mathcal D_0\); promoting them to
self-adjoint generator statements is a separate domain theorem.

\subsection{Exact Translation and Rotation Brackets}

Because \(h\) and \(p_a\) are multiplication operators,
\[
  i[h,p_a]=0.
\]
The canonical differential identity is
\[
  [X_a,p_b]=i\delta_{ab}.
\]
Therefore
\[
  i[P_c,J_{ab}]
  =
  \delta_{ac}P_b-\delta_{bc}P_a.
\]
This is independent of the dispersion.  It is kinematical: it comes from the
chosen momentum coordinate and the differential representation of rotations.

\subsection{Boost-Momentum Bracket}

For the boost candidate,
\[
\begin{aligned}
  [P_b,K_a]
  &=
  \frac{1}{2}\bigl([P_b,X_a]h+h[P_b,X_a]\bigr)\\
  &=
  -i\delta_{ab}h.
\end{aligned}
\]
Hence
\[
  i[P_b,K_a]=\delta_{ab}H .
\]
This part of the Poincare algebra is exact for every positive scalar
dispersion because it only uses \(K_a\) as the symmetrized first moment of
energy.

\subsection{Boost-Time Residual}

CA-24 derived
\[
  i[H,K_a]
  =
  \frac{1}{2}\partial_a\omega(k)^2.
\]
The continuum target with speed \(c=1\) wants \(P_a=k_a\).  The residual is
\[
  R_a(k)
  :=
  \frac{1}{2}\partial_a\omega(k)^2-k_a.
\]
For speed \(c\), replace \(k_a\) by \(c^2k_a\).  This is the first algebraic
filter on Hamiltonian coefficients.

\subsection{Rotation-Hamiltonian Residual}

The Hamiltonian commutes with rotations iff the dispersion is radial on the
patch:
\[
  i[h,J_{ab}]
  =
  k_b\partial_a\omega-k_a\partial_b\omega.
\]
Equivalently,
\[
  k_b\partial_a\omega^2-k_a\partial_b\omega^2=0
\]
where \(\omega>0\).  Hypercubic lattice symbols usually fail this globally on
the Brillouin zone but may satisfy it to leading order at a low-energy minimum.

\subsection{Boost-Boost Residual}

Write
\[
  B_a(k):=\frac{1}{2}\partial_a\omega(k)^2,
  \qquad
  R_a(k):=B_a(k)-k_a.
\]
Equivalently \(B_a=h\partial_a h\).  Since
\[
  K_a=i\left(h\partial_a+\frac{1}{2}\partial_ah\right),
\]
the commutator of the first-order differential parts gives
\[
  i[K_a,K_b]
  =
  J_{ab}+R_bX_a-R_aX_b.
\]
Here the possible divergence term vanishes because \(R\) is a gradient
residual:
\(\partial_aR_b=\partial_bR_a\).  Thus the boost-boost bracket closes to the
rotation bracket whenever the boost-time residual \(R_a\) vanishes on the patch.
The remaining caveat is analytic, not algebraic: all of these are unbounded
operators and still require a common self-adjointness/core theorem before this
becomes a continuum representation statement.

\subsection{Second Quantization}

On the finite-particle core of
\[
  \mathcal F_+(\mathfrak h_\epsilon),
\]
the many-body candidates are
\[
  H=d\Gamma(h),\quad
  P_a=d\Gamma(p_a),\quad
  J_{ab}=d\Gamma(J_{ab}^{(1)}),\quad
  K_a=d\Gamma(K_a^{(1)}),
\]
with a separate vacuum-energy convention if absolute Hamiltonian values are
used.  Commutators of second quantizations reproduce second quantizations of
the one-particle commutators on the finite-particle core.  Thus the compiler can
perform the first Lorentz tests at one-particle level and then lift the checked
residuals to the Gaussian Fock space.
